\begin{center}\begin{tikzpicture}[font=\small,>=stealth,line join=round]
\draw[gray!40,rounded corners=3pt] (0,0) rectangle (14.2,8.3);\node at (7.1,7.700000000000001) {例7：两条轴相距3，反射两次横移6};
\begin{scope}[shift={(2.242857142857143,3.7966666666666664)},x=0.5714285714285714cm,y=1.5333333333333332cm]\draw[->,gray] (-2,0)--(8.5,0) node[right]{$x$};\draw[->,gray] (0,-1.4)--(0,1.6) node[above]{$y$};\draw[blue,thick,solid,domain=-1.5:8,samples=180] plot(\x,{cos(deg(pi*(\x-1)/3))});\fill[orange!85!black] (0,0.5) circle(2.2pt);\fill[orange!85!black] (2,0.5) circle(2.2pt);\fill[orange!85!black] (6,0.5) circle(2.2pt);\foreach \a in {1,4}{\draw[violet,dashed] (\a,-1.3)--(\a,1.3);\node[above,violet] at (\a,1.3){$x=\a$};}\end{scope}\node[anchor=west,align=left,text width=5.0cm] at (8.4,4.0) {示例 $f=\cos\frac{\pi(x-1)}3$\\[9pt]轴 $x=1,\ x=4$\\[9pt]配对 $0\to2\to6$\\[9pt]三处高度都为 $\frac12$\\[9pt]$f(x+6)=f(x)$};
\node[text width=12.6cm,align=center] at (7.1,.65) {这条示例曲线帮助看反射；正文代数推导对所有满足题设的函数成立。};
\end{tikzpicture}\end{center}
